
\documentclass{article}
\usepackage{colortbl}
\usepackage{makecell}
\usepackage{multirow}
\usepackage{supertabular}

\begin{document}

\newcounter{utterance}

\centering \large Interaction Transcript for game `Functionigma Game', experiment `python\_easy', episode 1 with command{-}a{-}03{-}2025{-}t1.0.
\vspace{24pt}

{ \footnotesize  \setcounter{utterance}{1}
\setlength{\tabcolsep}{0pt}
\begin{supertabular}{c@{$\;$}|p{.15\linewidth}@{}p{.15\linewidth}p{.15\linewidth}p{.15\linewidth}p{.15\linewidth}p{.15\linewidth}}
    \# & \multicolumn{2}{c}{Player} && \multicolumn{2}{c}{Game Master} \\
    \hline

    \theutterance \stepcounter{utterance}  
    & & & \multicolumn{4}{p{0.6\linewidth}}{
        \cellcolor[rgb]{0.9,0.9,0.9}{
            \makecell[{{p{\linewidth}}}]{
                \texttt{\tiny{[P1$\langle$GM]}}
                \texttt{Your role:} \\
\texttt{You (the AI) will propose inputs for a hidden function so you can infer it.} \\
\\ 
\texttt{My role:} \\
\texttt{I (the user / GM) will return the output of a hidden function for each input you propose.} \\
\\ 
\texttt{Goal:} \\
\texttt{Infer the exact function:} \\
\texttt{(x: str) {-}> str} \\
\\ 
\texttt{Context:} \\
\texttt{Category: string.} \\
\texttt{This function takes 1 parameter(s) with names and types:} \\
\texttt{x: str} \\
\texttt{(Format shown below.)} \\
\\ 
\texttt{Important rules:} \\
\texttt{{-} You MUST provide inputs matching the declared parameter types.} \\
\texttt{  {-} For integers/floats: provide numeric literals (e.g. `{-}5`, `3.14`).} \\
\texttt{  {-} For strings: provide quoted strings (single or double quotes), e.g. `'hello'` or "hello".} \\
\texttt{  {-} For booleans: use `True` or `False`.} \\
\texttt{{-} Output values returned by the GM will be of the function's return type.} \\
\texttt{{-} Do not assume smoothness, monotonicity, etc. Test widely for edge cases.} \\
\\ 
\texttt{Protocol:} \\
\texttt{{-} Until you are certain, produce exactly one line **and nothing else**:} \\
\texttt{  `Input: v1, v2, v3, ...`} \\
\texttt{  where each vi has the exact type required (strings must be quoted).} \\
\texttt{{-} When certain, start your response with `GUESS:` followed by:} \\
\texttt{  1) The full guessed function as valid Python inside a triple{-}backtick fenced block labeled `python`.} \\
\texttt{  2) A short explanation of how you derived it.} \\
\\ 
\texttt{Examples:} \\
\texttt{{-} For a one{-}arg string function:} \\
\texttt{  `Input: 'hello'`} \\
\texttt{{-} For a two{-}arg int function:} \\
\texttt{  `Input: 3, {-}7`} \\
\\ 
\texttt{Formatting rules (strict):} \\
\texttt{{-} Commas between arguments.} \\
\texttt{{-} Strings must be quoted.} \\
\texttt{{-} No extra text, no explanation lines.} \\
\texttt{{-} If you violate the format, you will be prompted to resend.} \\
\\ 
\texttt{Begin. Generate your first input.} \\
            }
        }
    }
     \\ \\

    \theutterance \stepcounter{utterance}  
    & \multicolumn{4}{p{0.6\linewidth}}{
        \cellcolor[rgb]{0.561,0.737,0.561}{
            \makecell[{{p{\linewidth}}}]{
                \texttt{\tiny{[P1$\rangle$GM]}}
                \texttt{Input: 'a'} \\
            }
        }
    }
    & & \\ \\

    \theutterance \stepcounter{utterance}  
    & & & \multicolumn{4}{p{0.6\linewidth}}{
        \cellcolor[rgb]{0.9,0.9,0.9}{
            \makecell[{{p{\linewidth}}}]{
                \texttt{\tiny{[P1$\langle$GM]}}
                \texttt{Function Output: a} \\
            }
        }
    }
     \\ \\

    \theutterance \stepcounter{utterance}  
    & \multicolumn{4}{p{0.6\linewidth}}{
        \cellcolor[rgb]{0.561,0.737,0.561}{
            \makecell[{{p{\linewidth}}}]{
                \texttt{\tiny{[P1$\rangle$GM]}}
                \texttt{Input: 'abc'} \\
            }
        }
    }
    & & \\ \\

    \theutterance \stepcounter{utterance}  
    & & & \multicolumn{4}{p{0.6\linewidth}}{
        \cellcolor[rgb]{0.9,0.9,0.9}{
            \makecell[{{p{\linewidth}}}]{
                \texttt{\tiny{[P1$\langle$GM]}}
                \texttt{Function Output: cba} \\
            }
        }
    }
     \\ \\

    \theutterance \stepcounter{utterance}  
    & \multicolumn{4}{p{0.6\linewidth}}{
        \cellcolor[rgb]{0.561,0.737,0.561}{
            \makecell[{{p{\linewidth}}}]{
                \texttt{\tiny{[P1$\rangle$GM]}}
                \texttt{Input: 'hello'} \\
            }
        }
    }
    & & \\ \\

    \theutterance \stepcounter{utterance}  
    & & & \multicolumn{4}{p{0.6\linewidth}}{
        \cellcolor[rgb]{0.9,0.9,0.9}{
            \makecell[{{p{\linewidth}}}]{
                \texttt{\tiny{[P1$\langle$GM]}}
                \texttt{Function Output: olleh} \\
            }
        }
    }
     \\ \\

    \theutterance \stepcounter{utterance}  
    & \multicolumn{4}{p{0.6\linewidth}}{
        \cellcolor[rgb]{0.561,0.737,0.561}{
            \makecell[{{p{\linewidth}}}]{
                \texttt{\tiny{[P1$\rangle$GM]}}
                \texttt{Input: 'world'} \\
            }
        }
    }
    & & \\ \\

    \theutterance \stepcounter{utterance}  
    & & & \multicolumn{4}{p{0.6\linewidth}}{
        \cellcolor[rgb]{0.9,0.9,0.9}{
            \makecell[{{p{\linewidth}}}]{
                \texttt{\tiny{[P1$\langle$GM]}}
                \texttt{Function Output: dlrow} \\
            }
        }
    }
     \\ \\

    \theutterance \stepcounter{utterance}  
    & \multicolumn{4}{p{0.6\linewidth}}{
        \cellcolor[rgb]{0.561,0.737,0.561}{
            \makecell[{{p{\linewidth}}}]{
                \texttt{\tiny{[P1$\rangle$GM]}}
                \texttt{GUESS:} \\
\texttt{```python} \\
\texttt{def reverse\_string(x: str) {-}> str:} \\
\texttt{    return x{[}::{-}1{]}} \\
\texttt{```} \\
\\ 
\texttt{The function appears to reverse the input string. I tested with various strings, and the output was always the reversed version of the input. For example, `'a'` returns `'a'`, `'abc'` returns `'cba'`, `'hello'` returns `'olleh'`, and `'world'` returns `'dlrow'`. This pattern consistently indicates that the function is reversing the input string.} \\
            }
        }
    }
    & & \\ \\

    \theutterance \stepcounter{utterance}  
    & & & \multicolumn{2}{p{0.3\linewidth}}{
        \cellcolor[rgb]{0.9,0.9,0.9}{
            \makecell[{{p{\linewidth}}}]{
                \texttt{\tiny{[GM$|$GM]}}
                1.0
            }
        }
    }
    & & \\ \\

    \theutterance \stepcounter{utterance}  
    & & & \multicolumn{2}{p{0.3\linewidth}}{
        \cellcolor[rgb]{0.9,0.9,0.9}{
            \makecell[{{p{\linewidth}}}]{
                \texttt{\tiny{[GM$|$GM]}}
                \texttt{win} \\
            }
        }
    }
    & & \\ \\

    \theutterance \stepcounter{utterance}  
    & & & \multicolumn{2}{p{0.3\linewidth}}{
        \cellcolor[rgb]{0.9,0.9,0.9}{
            \makecell[{{p{\linewidth}}}]{
                \texttt{\tiny{[GM$|$GM]}}
                \texttt{That is correct! } \\
\texttt{The hidden function was:} \\
\\ 
\texttt{def rev\_str(x: str) {-}> str:} \\
\texttt{    return x{[}::{-}1{]}} \\
            }
        }
    }
    & & \\ \\

    \theutterance \stepcounter{utterance}  
    & & & \multicolumn{2}{p{0.3\linewidth}}{
        \cellcolor[rgb]{0.9,0.9,0.9}{
            \makecell[{{p{\linewidth}}}]{
                \texttt{\tiny{[GM$|$GM]}}
                \texttt{def reverse\_string(x: str) {-}> str:} \\
\texttt{    return x{[}::{-}1{]}} \\
            }
        }
    }
    & & \\ \\

\end{supertabular}
}

\end{document}
